%% kapa: exam=1 round=31 year=2020 period=2 form=A subject=appraisal_law subject_ko=감정평가관계법규 seq=36 number=36
%% source: Q-Net 감정평가사 기출문제 (cst00309) artlSeq=5207219 file=extracted/1008/5207219/A형/A형/2020년도 제31회 감정평가사 1차 시험_2교시_A형.pdf
%% conversion: kordoc 4.24.51 -> tools/gen1.py (auto); review=proofread-sonnet; flags=none
\begin{question}{36}
동산ㆍ채권 등의 담보에 관한 법령상 동산담보권에 관한 설명으로 옳은 것은?
\begin{choices}
  \item 동산담보권의 효력은 설정행위에 다른 약정이 있더라도 담보목적물에 부합된 물건과 종물(從物)에 미친다.
  \item 동산담보권은 피담보채권과 분리하여 타인에게 양도할 수 있다.
  \item 담보권설정자가 담보목적물을 점유하는 경우에 경매절차는 압류에 의하여 개시한다.
  \item 채무자의 변제를 원인으로 동산담보권의 실행을 중지함으로써 담보권자에게 손해가 발생하더라도 채무자가 그 손해를 배상하여야 하는 것은 아니다.
  \item 담보권자는 자기의 채권을 변제받기 위하여 담보목적물의 경매를 청구할 수 없다.
\end{choices}
\end{question}
\answer{3}
